Tenim dues molles idèntiques. Un objecte A de 100~g que penja d'una de les molles oscil·la amb un període d'1,00~s i amb una amplitud de 5,00~cm.

\begin{apartat}{1}
Volem que l'altra molla oscil·li amb la mateixa amplitud, però amb una freqüència doble que la de la molla de què penja l'objecte A. Quina massa hem de penjar a la segona molla?
\end{apartat}

\begin{apartat}{1}
Els dos objectes es deixen anar des de l'extrem inferior de l'oscil·lació. Representeu en una gràfica velocitat-temps la velocitat de cadascun dels objectes quan oscil·len durant 2~s en les condicions descrites. En la gràfica heu d'indicar clarament les escales dels eixos, les magnituds i les unitats. Durant els 2~s representats en la gràfica, en quins moments la diferència de fase entre els dos objectes és de $\pi$ radians?
\end{apartat}
